\scanpage{15}
\deleted{On dit}\inserted{Soit} $\bar T$ un voisinage
\deleted{tubulaire}\inserted{conique} de $Y\subset X$
dans $X$, posons
\[
(11)\qquad X^*=X\setminus\mathring T=(X\setminus\bar T)\cup\dot T
\]
c'est une partie fermée de $X$. \deleted{On dit que}

\uline{Prop.} Conditions équivalentes (où $\lambda\in]0,1[$ donné)

a) $\dot T$ est une partie-bord de $X^*=X\setminus\mathring T$

b) Il existe un homéomorphisme de $X$
avec lui-$=$, transformant $\bar T,\dot T$ en
$\deleted{\unsure{\lambda}}(\bar T_\lambda,\dot T_\lambda)$

b') Comme b, mais en supposant de plus
que l'homéomorphisme induit l'identité
sur $\bar T_{\lambda'}$ (où \deleted{$0<\lambda'<\lambda$}
$\lambda'$ est donné, $0<\lambda'<\lambda$).

c) Il existe un voisinage \deleted{tubulaire} conique $\bar T'$
de $Y$ dans $X$, \deleted{tel que et \unsure{a}} tel que $\bar T=\lambda\bar T'$,
l'isomorphisme étant compatible avec les
structures de voisinages coniques.

On dit alors que $\bar T$ est un vois. \deleted{tub.} conique
\uline{strict} dans $X$.

\margin{Il vaut mieux de définir les voisinages coniques
stricts comme \unsure{des} \uline{germes} de voisinage
plus grand $(1+\varepsilon)\bar T\unsure{-R_{>0}}$
$(\varepsilon>0)$}

\uline{Corollaire} Quand $\bar T$ est vois. \deleted{tub.} con. strict dans $X$,
$\dot T$ est la frontière de $\bar T$ dans $X$.

\uline{Corollaire} Si \deleted{il est}\inserted{$\bar T$} voisinage
\deleted{tubulaire}\inserted{conique} de $Y$
dans $X$, les $\bar T_\lambda$ $(\lambda\in]0,1[)$ sont des voisinages
\deleted{tubulaires}\inserted{coniques} stricts.
\margin{14}
